\documentclass[a4paper,12pt]{article}
\usepackage[utf8]{inputenc}
\usepackage[T2A]{fontenc}
\usepackage[russian]{babel}
\usepackage{amsmath, amssymb, amsthm}
\usepackage{fullpage}
\usepackage{enumitem}

\newtheorem{theorem}{Теорема}[section]
\newtheorem{lemma}[theorem]{Лемма}
\title{}
\author{}
\date{}

\begin{document}
	\maketitle
	
	\section*{Основная часть}


	
	\subsection*{Задача 1: Выпуклость функции}
	
	Рассмотрим функцию $f: \mathbb{R}^2 \rightarrow \mathbb{R}$, заданную как $f(x) = f(x_1, x_2) = x_1^2 x_2^2$. Для определения выпуклости функции, мы можем исследовать ее гессиан.
	
	Сначала найдем первые частные производные:
	\begin{align*}
		\frac{\partial f}{\partial x_1} &= 2x_1 x_2^2 \\
		\frac{\partial f}{\partial x_2} &= 2x_1^2 x_2
	\end{align*}
	
	Теперь найдем вторые частные производные:
	\begin{align*}
		\frac{\partial^2 f}{\partial x_1^2} &= \frac{\partial}{\partial x_1} (2x_1 x_2^2) = 2x_2^2 \\
		\frac{\partial^2 f}{\partial x_2^2} &= \frac{\partial}{\partial x_2} (2x_1^2 x_2) = 2x_1^2 \\
		\frac{\partial^2 f}{\partial x_1 \partial x_2} &= \frac{\partial}{\partial x_1} (2x_1^2 x_2) = 4x_1 x_2 \\
		\frac{\partial^2 f}{\partial x_2 \partial x_1} &= \frac{\partial}{\partial x_2} (2x_1 x_2^2) = 4x_1 x_2
	\end{align*}
	
	Гессиан функции $f$ имеет вид:
	$$
	H_f(x) = \begin{pmatrix}
		\frac{\partial^2 f}{\partial x_1^2} & \frac{\partial^2 f}{\partial x_1 \partial x_2} \\
		\frac{\partial^2 f}{\partial x_2 \partial x_1} & \frac{\partial^2 f}{\partial x_2^2}
	\end{pmatrix} =
	\begin{pmatrix}
		2x_2^2 & 4x_1 x_2 \\
		4x_1 x_2 & 2x_1^2
	\end{pmatrix}
	$$
	
	Для того чтобы функция $f$ была выпуклой, ее гессиан должен быть положительно полуопределенным для всех $x \in \mathbb{R}^2$. Для матрицы $2 \times 2$, это означает, что главные миноры должны быть неотрицательными.
	
	Первый главный минор:
	$$
	\Delta_1 = 2x_2^2 \ge 0 \quad \forall x_2 \in \mathbb{R}
	$$
	
	Второй главный минор (определитель гессиана):
	$$
	\Delta_2 = \det(H_f(x)) = (2x_2^2)(2x_1^2) - (4x_1 x_2)(4x_1 x_2) = 4x_1^2 x_2^2 - 16x_1^2 x_2^2 = -12x_1^2 x_2^2
	$$
	
	Для положительной полуопределенности гессиана, необходимо, чтобы $\Delta_2 \ge 0$. Однако, $-12x_1^2 x_2^2 \ge 0$ только в том случае, если $x_1 = 0$ или $x_2 = 0$. Например, если $x_1 = 1$ и $x_2 = 1$, то $\Delta_2 = -12 < 0$.
	
	Так как определитель гессиана может быть отрицательным, функция $f(x) = x_1^2 x_2^2$ не является выпуклой на $\mathbb{R}^2$.
	
	
	
	\subsection*{Задача 2}

	\textbf{Решение.} Для каждой координаты функция $x\mapsto x^4$ имеет вторую производную
	\[
	\frac{d^2}{dx^2} x^4 = 12x^2.
	\]
	Так как $12x^2\ge 0$ для всех $x\in\mathbb{R}$, функция $x^4$ является выпуклой. Следовательно, сумма таких функций по разным координатам также выпукла.
	
	Для $\mu$-сильной выпуклости требуется, чтобы для всех $x\in\mathbb{R}^d$ гессиан $H_f(x)$ удовлетворял неравенству
	\[
	H_f(x) \succeq \mu I.
	\]
	При вычислении гессиана получаем диагональную матрицу с элементами $12x_i^2$. Заметим, что при $x=0$ все диагональные элементы равны нулю, то есть $H_f(0)=0$. Таким образом, не существует $\mu>0$, для которого $H_f(x)\succeq \mu I$ для всех $x$. 
	
	Ответ: функция $f$ выпукла, но не является $\mu$-сильно выпуклой ни для какого $\mu>0$.
	
	\vspace{0.5em}
	
	\subsection*{Задача 3}

	
	\textbf{Решение.} Из классических результатов выпуклого анализа и теории спектральных функций известно следующее:
	\begin{enumerate}[label=1)]
		\item Функция, сопоставляющая симметричной матрице её наибольшее собственное значение, является выпуклой.
		\item Функция, сопоставляющая симметричной матрице её наименьшее собственное значение, является вогнутой.
	\end{enumerate}
	Таким образом, ответы: 
	\begin{itemize}
		\item $f(X)=\lambda_{\max}(X)$ --- выпуклая;
		\item $f(X)=\lambda_{\min}(X)$ --- вогнутая.
	\end{itemize}
	
	\vspace{0.5em}
	
	\subsection*{Задача 4}
	\textbf{Условие.} Пусть $f:\operatorname{dom}f\to\mathbb{R}$ --- функция с областью определения $\operatorname{dom}f\subset\mathbb{R}^d$. Докажите, что $f$ выпукла тогда и только тогда, когда её сужение на любую прямую является выпуклым. Формально: для любых $x_0\in\operatorname{dom}f$ и $u\in\mathbb{R}^d$ функция 
	\[
	g(t)=f(x_0+tu),\quad t\in\{t\in\mathbb{R}:x_0+tu\in\operatorname{dom}f\},
	\]
	является выпуклой.
	
	\textbf{Доказательство.}
	
	\emph{(Необходимость).} Пусть $f$ выпукла, т.е. для любых $x,y\in\operatorname{dom}f$ и $\lambda\in[0,1]$ выполнено
	\[
	f(\lambda x+(1-\lambda)y)\le\lambda f(x)+(1-\lambda)f(y).
	\]
	Выберем произвольное $x_0\in\operatorname{dom}f$ и $u\in\mathbb{R}^d$. Для любых $t_1,t_2$ из области определения функции $g$ и для любого $\lambda\in[0,1]$ рассмотрим точки $x_0+t_1u$ и $x_0+t_2u$. Тогда
	\[
	\begin{aligned}
		g(\lambda t_1+(1-\lambda)t_2)&=f\bigl(x_0+(\lambda t_1+(1-\lambda)t_2)u\bigr)\\[1mm]
		&=f\Bigl(\lambda (x_0+t_1u)+(1-\lambda)(x_0+t_2u)\Bigr)\\[1mm]
		&\le \lambda f(x_0+t_1u)+(1-\lambda) f(x_0+t_2u)\\[1mm]
		&=\lambda g(t_1)+(1-\lambda)g(t_2).
	\end{aligned}
	\]
	Таким образом, $g$ выпукла.
	
	\emph{(Достаточность).} Пусть для любого $x_0\in\operatorname{dom}f$ и любого направления $u\in\mathbb{R}^d$ функция 
	\[
	g(t)=f(x_0+tu)
	\]
	выпукла на своем множестве определения. Возьмём произвольные $x,y\in\operatorname{dom}f$ и положим $u=y-x$; тогда можно записать
	\[
	x=\;x,\quad y=x+1\cdot u.
	\]
	Определим функцию $g(t)=f(x+tu)$, которая по предположению выпукла для $t$ таких, что $x+tu\in\operatorname{dom}f$. При $t=0$ и $t=1$ имеем $g(0)=f(x)$ и $g(1)=f(y)$, а выпуклость $g$ даёт
	\[
	f(\lambda x+(1-\lambda)y)=g(1-\lambda)\le (1-\lambda)g(0)+\lambda g(1)=(1-\lambda)f(x)+\lambda f(y)
	\]
	для всех $\lambda\in[0,1]$. Это и есть определение выпуклости функции $f$. 
	
	Таким образом, $f$ выпукла тогда и только тогда, когда её сужение на любую прямую выпукло.
	
	\vspace{1em}
	
	\section*{Дополнительная часть}
	
	\subsection*{Задача 1}
	\textbf{Решение.} Данное неравенство является неотрицательностью дивергенции Кульбака–Лейблера:
	\[
	D_{\mathrm{KL}}(p\|q)=\sum_{i=1}^d p_i \ln\frac{p_i}{q_i}\ge 0,
	\]
	при этом равенство достигается тогда и только тогда, когда $p_i=q_i$ для всех $i$. Это утверждение можно доказать с помощью неравенства Йенсена, применив его к выпуклой функции $\varphi(x)=x\ln x$.
	
	\vspace{0.5em}
	
	\subsection*{Задача 2}
	\textbf{Решение.} Рассмотрим функцию 
	\[
	F(x)=\int_0^x g(t)\,dt,
	\]
	таким образом $f(x)=\frac{F(x)}{x}$. Заметим, что при $g(0)=0$ функция $F$ удовлетворяет $F(0)=0$. В таких случаях можно доказать, что функция $f(x)$ является неубывающей, а затем, анализируя её вторую производную, установить выпуклость. 
	
	Альтернативный подход --- рассмотреть неравенство Чебышева для интегралов или использовать следующее соображение: для $0 < x_1 < x_2$ имеет место
	\[
	\frac{F(x_2)}{x_2} - \frac{F(x_1)}{x_1} \ge 0,
	\]
	что эквивалентно неубыванию функции $f(x)$. Далее, достаточно проверить, что производная $f'(x)$ неубывает. Детальный расчёт второго производного (с использованием формулы Ферма для дифференцирования отношения) приводит к тому, что при выпуклости $g$ выполняется $f''(x)\ge 0$, откуда и следует выпуклость $f$.
	
	\vspace{0.5em}
	
	\subsection*{Задача 3}
	\textbf{Решение.} Из известных результатов матричного анализа следует, что функция, сопоставляющая матрице $X\in S_d^{++}$ значение $\operatorname{Tr}(X^{-1})$, является выпуклой. Это можно доказать, используя свойства обратной матрицы и соответствующую матричную неравномерность (например, через второй порядковый разложение или операторную выпуклость функции $X\mapsto X^{-1}$).
	
	Ответ: функция $f(X)=\operatorname{Tr}(X^{-1})$ выпукла на $S_d^{++}$.
	
	\vspace{0.5em}
	
	\subsection*{Задача 4}

	
	\textbf{Решение.}
	
	\textbf{Шаг 1.} Для набора положительных чисел $\{a_i\}$ и весов $\{\alpha_i\}$, где $\alpha_i\ge 0$, $\sum_{i=1}^d \alpha_i=1$, неравенство Йенсена гласит:
	\[
	-\ln\Biggl(\sum_{i=1}^d \alpha_i a_i\Biggr) \le \sum_{i=1}^d \alpha_i (-\ln a_i).
	\]
	Эквивалентно,
	\[
	\ln\Biggl(\sum_{i=1}^d \alpha_i a_i\Biggr) \ge \sum_{i=1}^d \alpha_i \ln a_i.
	\]
	
	\textbf{Шаг 2.} Выберем
	\[
	\alpha_i = \frac{|x_i|^p}{\sum_{j=1}^d |x_j|^p},
	\]
	а также положим
	\[
	a_i = \frac{|y_i|}{|x_i|^{p/q}},
	\]
	при условии, что все $x_i\neq 0$ (остальные случаи разбиваются аналогичным рассуждением, а предел можно получить переходом к пределу).
	
	Тогда неравенство Йенсена приводит к соотношению, которое после преобразований (учитывая соотношение $\frac{1}{p}+\frac{1}{q}=1$) эквивалентно неравенству Гёльдера.
	
	\textbf{Комментарий.} Детальное преобразование включает логарифмирование, применение экспоненты, а затем использование определения сопряжённых норм. В итоге получаем:
	\[
	\sum_{i=1}^d x_i y_i \le \Biggl(\sum_{i=1}^d |x_i|^p\Biggr)^{1/p}\Biggl(\sum_{i=1}^d |y_i|^q\Biggr)^{1/q}.
	\]
	
	Таким образом, неравенство Гёльдера доказано с использованием неравенства Йенсена для функции $-\ln x$.
	
	\vspace{1em}
	
	
\end{document}
